% Figure: donating regions of one x sweep (Chapter 5). Cells, PLIC and solid
% lines, donating rectangles, their open and liquid parts and the face
% velocities: generated/vof_donating.tex (vof.jl: shear fixture for panels a
% and b, cut fixture for panel c; regions from the independent evaluation that
% reproduces the code's fluxes). Labels hand-written.
\begin{figure}
  \centering
  \def\fxvnone{%
    \def\fxvcell##1##2##3##4##5##6##7{}\def\fxvliquid##1{}\def\fxvplic##1##2##3##4{}%
    \def\fxvsolid##1{}\def\fxvwall##1##2##3##4{}\def\fxvdr##1##2##3##4{}%
    \def\fxvvel##1##2##3##4{}}%
  % Passes: 1 liquid and solid, 2 swept open and liquid parts, 3 lines, 4
  % donating rectangles, 5 frozen indicator. #1: panel number.
  \def\fxvpasses#1{%
    \def\fxvpanel{#1}%
    \fxvnone
    \def\fxvliquid##1{\fill[tr liquid] ##1;}%
    \def\fxvsolid##1{\fill[white] ##1; \fill[tr solid] ##1;}%
    % Generated by docs/techref/figures/generate.jl; do not edit.
\ifnum\fxvpanel=1
\fxvcell{0}{3}{1}{4}{0.00}{0}{0}
\fxvcell{1}{3}{2}{4}{0.01}{0}{0}
\fxvliquid{(2,3.811366) -- (2,4) -- (1.929225,4) -- cycle}
\fxvplic{2}{3.811366}{1.929225}{4}
\fxvcell{2}{3}{3}{4}{0.92}{1}{0}
\fxvliquid{(2.333083,3) -- (3,3) -- (3,4) -- (2,4) -- (2,3.50418) -- cycle}
\fxvplic{2.333083}{3}{2}{3.50418}
\fxvcell{0}{4}{1}{5}{0.00}{0}{0}
\fxvcell{1}{4}{2}{5}{0.03}{0}{0}
\fxvliquid{(1.912944,4) -- (2,4) -- (2,4.58673) -- cycle}
\fxvplic{1.912944}{4}{2}{4.58673}
\fxvcell{2}{4}{3}{5}{0.99}{1}{0}
\fxvliquid{(2,4) -- (3,4) -- (3,5) -- (2.100546,5) -- (2,4.739856) -- cycle}
\fxvplic{2.100546}{5}{2}{4.739856}
\fxvcell{0}{5}{1}{6}{0.00}{0}{0}
\fxvcell{1}{5}{2}{6}{0.00}{0}{0}
\fxvcell{2}{5}{3}{6}{0.63}{1}{0}
\fxvliquid{(2,5) -- (3,5) -- (3,6) -- (2.955594,6) -- (2,5.217152) -- cycle}
\fxvplic{2.955594}{6}{2}{5.217152}
\fxvdr{(0.784912,3) -- (1,3) -- (1,4) -- (0.784912,4) -- cycle}{(0.784912,3) -- (1,3) -- (1,4) -- (0.784912,4) -- cycle}{}{0.22}
\fxvdr{(0.828196,4) -- (1,4) -- (1,5) -- (0.828196,5) -- cycle}{(0.828196,4) -- (1,4) -- (1,5) -- (0.828196,5) -- cycle}{}{0.17}
\fxvdr{(0.858803,5) -- (1,5) -- (1,6) -- (0.858803,6) -- cycle}{(0.858803,5) -- (1,5) -- (1,6) -- (0.858803,6) -- cycle}{}{0.14}
\fxvdr{(1.784912,3) -- (2,3) -- (2,4) -- (1.784912,4) -- cycle}{(1.784912,3) -- (2,3) -- (2,4) -- (1.784912,4) -- cycle}{(2,3.811366) -- (2,4) -- (1.929225,4) -- cycle}{0.22}
\fxvdr{(1.828196,4) -- (2,4) -- (2,5) -- (1.828196,5) -- cycle}{(1.828196,4) -- (2,4) -- (2,5) -- (1.828196,5) -- cycle}{(1.912944,4) -- (2,4) -- (2,4.58673) -- cycle}{0.17}
\fxvdr{(1.858803,5) -- (2,5) -- (2,6) -- (1.858803,6) -- cycle}{(1.858803,5) -- (2,5) -- (2,6) -- (1.858803,6) -- cycle}{}{0.14}
\fxvdr{(2.784912,3) -- (3,3) -- (3,4) -- (2.784912,4) -- cycle}{(2.784912,3) -- (3,3) -- (3,4) -- (2.784912,4) -- cycle}{(2.784912,3) -- (3,3) -- (3,4) -- (2.784912,4) -- cycle}{0.22}
\fxvdr{(2.828196,4) -- (3,4) -- (3,5) -- (2.828196,5) -- cycle}{(2.828196,4) -- (3,4) -- (3,5) -- (2.828196,5) -- cycle}{(2.828196,4) -- (3,4) -- (3,5) -- (2.828196,5) -- cycle}{0.17}
\fxvdr{(2.858803,5) -- (3,5) -- (3,6) -- (2.858803,6) -- cycle}{(2.858803,5) -- (3,5) -- (3,6) -- (2.858803,6) -- cycle}{(2.858803,5) -- (3,5) -- (3,6) -- (2.955594,6) -- (2.858803,5.920706) -- cycle}{0.14}
\fxvvel{1}{0}{3.5}{1.112}
\fxvvel{1}{0}{4.5}{0.888}
\fxvvel{1}{0}{5.5}{0.730}
\fxvvel{1}{1}{3.5}{1.112}
\fxvvel{1}{1}{4.5}{0.888}
\fxvvel{1}{1}{5.5}{0.730}
\fxvvel{1}{2}{3.5}{1.112}
\fxvvel{1}{2}{4.5}{0.888}
\fxvvel{1}{2}{5.5}{0.730}
\fxvvel{1}{3}{3.5}{1.112}
\fxvvel{1}{3}{4.5}{0.888}
\fxvvel{1}{3}{5.5}{0.730}
\fi
\ifnum\fxvpanel=2
\fxvcell{3}{5}{4}{6}{1.00}{1}{0}
\fxvliquid{(3,5) -- (4,5) -- (4,6) -- (3,6) -- cycle}
\fxvcell{3}{6}{4}{7}{0.39}{0}{0}
\fxvliquid{(3,6) -- (4,6) -- (4,6.455724) -- (3,6.320332) -- cycle}
\fxvplic{4}{6.455724}{3}{6.320332}
\fxvcell{4}{5}{4.5}{5.5}{1.00}{1}{0}
\fxvliquid{(4,5) -- (4.5,5) -- (4.5,5.5) -- (4,5.5) -- cycle}
\fxvcell{4.5}{5}{5}{5.5}{1.00}{1}{0}
\fxvliquid{(4.5,5) -- (5,5) -- (5,5.5) -- (4.5,5.5) -- cycle}
\fxvcell{4}{5.5}{4.5}{6}{1.00}{1}{0}
\fxvliquid{(4,5.5) -- (4.5,5.5) -- (4.5,6) -- (4,6) -- cycle}
\fxvcell{4.5}{5.5}{5}{6}{1.00}{1}{0}
\fxvliquid{(4.5,5.5) -- (5,5.5) -- (5,6) -- (4.5,6) -- cycle}
\fxvcell{4}{6}{4.5}{6.5}{1.00}{1}{0}
\fxvliquid{(4,6) -- (4.5,6) -- (4.5,6.5) -- (4,6.5) -- cycle}
\fxvcell{4.5}{6}{5}{6.5}{0.98}{1}{0}
\fxvliquid{(4.5,6) -- (5,6) -- (5,6.451796) -- (4.760683,6.5) -- (4.5,6.5) -- cycle}
\fxvplic{5}{6.451796}{4.760683}{6.5}
\fxvcell{4}{6.5}{4.5}{7}{0.09}{0}{0}
\fxvliquid{(4,6.5) -- (4.5,6.5) -- (4.5,6.565009) -- (4,6.52504) -- cycle}
\fxvplic{4.5}{6.565009}{4}{6.52504}
\fxvcell{4.5}{6.5}{5}{7}{0.03}{0}{0}
\fxvliquid{(4.5,6.5) -- (4.757451,6.5) -- (4.5,6.551857) -- cycle}
\fxvplic{4.757451}{6.5}{4.5}{6.551857}
\fxvdr{(3.854498,5) -- (4,5) -- (4,5.5) -- (3.854498,5.5) -- cycle}{(3.854498,5) -- (4,5) -- (4,5.5) -- (3.854498,5.5) -- cycle}{(3.854498,5) -- (4,5) -- (4,5.5) -- (3.854498,5.5) -- cycle}{0.15}
\fxvdr{(3.863107,5.5) -- (4,5.5) -- (4,6) -- (3.863107,6) -- cycle}{(3.863107,5.5) -- (4,5.5) -- (4,6) -- (3.863107,6) -- cycle}{(3.863107,5.5) -- (4,5.5) -- (4,6) -- (3.863107,6) -- cycle}{0.14}
\fxvdr{(3.863107,6) -- (4,6) -- (4,6.5) -- (3.863107,6.5) -- cycle}{(3.863107,6) -- (4,6) -- (4,6.5) -- (3.863107,6.5) -- cycle}{(3.863107,6) -- (4,6) -- (4,6.455724) -- (3.863107,6.43719) -- cycle}{0.14}
\fxvdr{(3.854498,6.5) -- (4,6.5) -- (4,7) -- (3.854498,7) -- cycle}{(3.854498,6.5) -- (4,6.5) -- (4,7) -- (3.854498,7) -- cycle}{}{0.15}
\fxvdr{(4.354498,5) -- (4.5,5) -- (4.5,5.5) -- (4.354498,5.5) -- cycle}{(4.354498,5) -- (4.5,5) -- (4.5,5.5) -- (4.354498,5.5) -- cycle}{(4.354498,5) -- (4.5,5) -- (4.5,5.5) -- (4.354498,5.5) -- cycle}{0.29}
\fxvdr{(4.363107,5.5) -- (4.5,5.5) -- (4.5,6) -- (4.363107,6) -- cycle}{(4.363107,5.5) -- (4.5,5.5) -- (4.5,6) -- (4.363107,6) -- cycle}{(4.363107,5.5) -- (4.5,5.5) -- (4.5,6) -- (4.363107,6) -- cycle}{0.27}
\fxvdr{(4.363107,6) -- (4.5,6) -- (4.5,6.5) -- (4.363107,6.5) -- cycle}{(4.363107,6) -- (4.5,6) -- (4.5,6.5) -- (4.363107,6.5) -- cycle}{(4.363107,6) -- (4.5,6) -- (4.5,6.5) -- (4.363107,6.5) -- cycle}{0.27}
\fxvdr{(4.354498,6.5) -- (4.5,6.5) -- (4.5,7) -- (4.354498,7) -- cycle}{(4.354498,6.5) -- (4.5,6.5) -- (4.5,7) -- (4.354498,7) -- cycle}{(4.354498,6.5) -- (4.5,6.5) -- (4.5,6.565009) -- (4.354498,6.553378) -- cycle}{0.29}
\fxvdr{(4.854498,5) -- (5,5) -- (5,5.5) -- (4.854498,5.5) -- cycle}{(4.854498,5) -- (5,5) -- (5,5.5) -- (4.854498,5.5) -- cycle}{(4.854498,5) -- (5,5) -- (5,5.5) -- (4.854498,5.5) -- cycle}{0.29}
\fxvdr{(4.863107,5.5) -- (5,5.5) -- (5,6) -- (4.863107,6) -- cycle}{(4.863107,5.5) -- (5,5.5) -- (5,6) -- (4.863107,6) -- cycle}{(4.863107,5.5) -- (5,5.5) -- (5,6) -- (4.863107,6) -- cycle}{0.27}
\fxvdr{(4.863107,6) -- (5,6) -- (5,6.5) -- (4.863107,6.5) -- cycle}{(4.863107,6) -- (5,6) -- (5,6.5) -- (4.863107,6.5) -- cycle}{(4.863107,6) -- (5,6) -- (5,6.451796) -- (4.863107,6.479369) -- cycle}{0.27}
\fxvdr{(4.854498,6.5) -- (5,6.5) -- (5,7) -- (4.854498,7) -- cycle}{(4.854498,6.5) -- (5,6.5) -- (5,7) -- (4.854498,7) -- cycle}{}{0.29}
\fxvvel{1}{3}{5.5}{0.730}
\fxvvel{1}{3}{6.5}{0.730}
\fxvvel{1}{4}{5.25}{0.752}
\fxvvel{1}{4}{5.75}{0.708}
\fxvvel{1}{4}{6.25}{0.708}
\fxvvel{1}{4}{6.75}{0.752}
\fxvvel{1}{4.5}{5.25}{0.752}
\fxvvel{1}{4.5}{5.75}{0.708}
\fxvvel{1}{4.5}{6.25}{0.708}
\fxvvel{1}{4.5}{6.75}{0.752}
\fxvvel{1}{5}{5.25}{0.752}
\fxvvel{1}{5}{5.75}{0.708}
\fxvvel{1}{5}{6.25}{0.708}
\fxvvel{1}{5}{6.75}{0.752}
\fi
\ifnum\fxvpanel=3
\fxvcell{1}{0}{2}{1}{1.00}{1}{1}
\fxvsolid{(1,0) -- (2,0) -- (2,0.45) -- (1,0.25) -- cycle}
\fxvwall{2}{0.45}{1}{0.25}
\fxvliquid{(2,0.45) -- (2,1) -- (1,1) -- (1,0.25) -- cycle}
\fxvcell{2}{0}{3}{1}{0.66}{1}{1}
\fxvsolid{(2,0) -- (3,0) -- (3,0.55) -- (2,0.35) -- cycle}
\fxvwall{3}{0.55}{2}{0.35}
\fxvliquid{(2.7,1) -- (2,1) -- (2,0.35) -- (2.537234,0.457447) -- cycle}
\fxvplic{2.7}{1}{2.537234}{0.457447}
\fxvcell{3}{0}{4}{1}{0.00}{0}{1}
\fxvsolid{(3,0) -- (4,0) -- (4,0.65) -- (3,0.45) -- cycle}
\fxvwall{4}{0.65}{3}{0.45}
\fxvcell{1}{1}{2}{2}{1.00}{1}{0}
\fxvliquid{(1,1) -- (2,1) -- (2,2) -- (1,2) -- cycle}
\fxvcell{2}{1}{3}{2}{0.85}{1}{0}
\fxvliquid{(2,1) -- (2.7,1) -- (3,2) -- (2,2) -- cycle}
\fxvplic{2.7}{1}{3}{2}
\fxvcell{3}{1}{4}{2}{0.00}{0}{0}
\fxvdr{(1.577856,0) -- (2,0) -- (2,1) -- (1.577856,1) -- cycle}{(2,0.45) -- (2,1) -- (1.577856,1) -- (1.577856,0.365571) -- cycle}{(2,0.45) -- (2,1) -- (1.577856,1) -- (1.577856,0.365571) -- cycle}{0.42}
\fxvdr{(1.583333,1) -- (2,1) -- (2,2) -- (1.583333,2) -- cycle}{(1.583333,1) -- (2,1) -- (2,2) -- (1.583333,2) -- cycle}{(1.583333,1) -- (2,1) -- (2,2) -- (1.583333,2) -- cycle}{0.42}
\fxvdr{(2.576831,0) -- (3,0) -- (3,1) -- (2.576831,1) -- cycle}{(3,0.55) -- (3,1) -- (2.576831,1) -- (2.576831,0.465366) -- cycle}{(2.7,1) -- (2.576831,1) -- (2.576831,0.589436) -- cycle}{0.42}
\fxvdr{(2.583333,1) -- (3,1) -- (3,2) -- (2.583333,2) -- cycle}{(2.583333,1) -- (3,1) -- (3,2) -- (2.583333,2) -- cycle}{(2.583333,1) -- (2.7,1) -- (3,2) -- (2.583333,2) -- cycle}{0.42}
\fxvdr{(3.575335,0) -- (4,0) -- (4,1) -- (3.575335,1) -- cycle}{(4,0.65) -- (4,1) -- (3.575335,1) -- (3.575335,0.565067) -- cycle}{}{0.42}
\fxvdr{(3.583333,1) -- (4,1) -- (4,2) -- (3.583333,2) -- cycle}{(3.583333,1) -- (4,1) -- (4,2) -- (3.583333,2) -- cycle}{}{0.42}
\fxvvel{1}{1}{0.5}{1.000}
\fxvvel{1}{1}{1.5}{1.000}
\fxvvel{1}{2}{0.5}{1.000}
\fxvvel{1}{2}{1.5}{1.000}
\fxvvel{1}{3}{0.5}{1.000}
\fxvvel{1}{3}{1.5}{1.000}
\fxvvel{1}{4}{0.5}{1.000}
\fxvvel{1}{4}{1.5}{1.000}
\fi
%
    \fxvnone
    \def\fxvdr##1##2##3##4{\def\fxvarg{##2}\ifx\fxvarg\empty\else\fill[tr dr open] ##2;\fi
      \def\fxvarg{##3}\ifx\fxvarg\empty\else\fill[tr liquid] ##3; \fill[tr flux] ##3;\fi}%
    % Generated by docs/techref/figures/generate.jl; do not edit.
\ifnum\fxvpanel=1
\fxvcell{0}{3}{1}{4}{0.00}{0}{0}
\fxvcell{1}{3}{2}{4}{0.01}{0}{0}
\fxvliquid{(2,3.811366) -- (2,4) -- (1.929225,4) -- cycle}
\fxvplic{2}{3.811366}{1.929225}{4}
\fxvcell{2}{3}{3}{4}{0.92}{1}{0}
\fxvliquid{(2.333083,3) -- (3,3) -- (3,4) -- (2,4) -- (2,3.50418) -- cycle}
\fxvplic{2.333083}{3}{2}{3.50418}
\fxvcell{0}{4}{1}{5}{0.00}{0}{0}
\fxvcell{1}{4}{2}{5}{0.03}{0}{0}
\fxvliquid{(1.912944,4) -- (2,4) -- (2,4.58673) -- cycle}
\fxvplic{1.912944}{4}{2}{4.58673}
\fxvcell{2}{4}{3}{5}{0.99}{1}{0}
\fxvliquid{(2,4) -- (3,4) -- (3,5) -- (2.100546,5) -- (2,4.739856) -- cycle}
\fxvplic{2.100546}{5}{2}{4.739856}
\fxvcell{0}{5}{1}{6}{0.00}{0}{0}
\fxvcell{1}{5}{2}{6}{0.00}{0}{0}
\fxvcell{2}{5}{3}{6}{0.63}{1}{0}
\fxvliquid{(2,5) -- (3,5) -- (3,6) -- (2.955594,6) -- (2,5.217152) -- cycle}
\fxvplic{2.955594}{6}{2}{5.217152}
\fxvdr{(0.784912,3) -- (1,3) -- (1,4) -- (0.784912,4) -- cycle}{(0.784912,3) -- (1,3) -- (1,4) -- (0.784912,4) -- cycle}{}{0.22}
\fxvdr{(0.828196,4) -- (1,4) -- (1,5) -- (0.828196,5) -- cycle}{(0.828196,4) -- (1,4) -- (1,5) -- (0.828196,5) -- cycle}{}{0.17}
\fxvdr{(0.858803,5) -- (1,5) -- (1,6) -- (0.858803,6) -- cycle}{(0.858803,5) -- (1,5) -- (1,6) -- (0.858803,6) -- cycle}{}{0.14}
\fxvdr{(1.784912,3) -- (2,3) -- (2,4) -- (1.784912,4) -- cycle}{(1.784912,3) -- (2,3) -- (2,4) -- (1.784912,4) -- cycle}{(2,3.811366) -- (2,4) -- (1.929225,4) -- cycle}{0.22}
\fxvdr{(1.828196,4) -- (2,4) -- (2,5) -- (1.828196,5) -- cycle}{(1.828196,4) -- (2,4) -- (2,5) -- (1.828196,5) -- cycle}{(1.912944,4) -- (2,4) -- (2,4.58673) -- cycle}{0.17}
\fxvdr{(1.858803,5) -- (2,5) -- (2,6) -- (1.858803,6) -- cycle}{(1.858803,5) -- (2,5) -- (2,6) -- (1.858803,6) -- cycle}{}{0.14}
\fxvdr{(2.784912,3) -- (3,3) -- (3,4) -- (2.784912,4) -- cycle}{(2.784912,3) -- (3,3) -- (3,4) -- (2.784912,4) -- cycle}{(2.784912,3) -- (3,3) -- (3,4) -- (2.784912,4) -- cycle}{0.22}
\fxvdr{(2.828196,4) -- (3,4) -- (3,5) -- (2.828196,5) -- cycle}{(2.828196,4) -- (3,4) -- (3,5) -- (2.828196,5) -- cycle}{(2.828196,4) -- (3,4) -- (3,5) -- (2.828196,5) -- cycle}{0.17}
\fxvdr{(2.858803,5) -- (3,5) -- (3,6) -- (2.858803,6) -- cycle}{(2.858803,5) -- (3,5) -- (3,6) -- (2.858803,6) -- cycle}{(2.858803,5) -- (3,5) -- (3,6) -- (2.955594,6) -- (2.858803,5.920706) -- cycle}{0.14}
\fxvvel{1}{0}{3.5}{1.112}
\fxvvel{1}{0}{4.5}{0.888}
\fxvvel{1}{0}{5.5}{0.730}
\fxvvel{1}{1}{3.5}{1.112}
\fxvvel{1}{1}{4.5}{0.888}
\fxvvel{1}{1}{5.5}{0.730}
\fxvvel{1}{2}{3.5}{1.112}
\fxvvel{1}{2}{4.5}{0.888}
\fxvvel{1}{2}{5.5}{0.730}
\fxvvel{1}{3}{3.5}{1.112}
\fxvvel{1}{3}{4.5}{0.888}
\fxvvel{1}{3}{5.5}{0.730}
\fi
\ifnum\fxvpanel=2
\fxvcell{3}{5}{4}{6}{1.00}{1}{0}
\fxvliquid{(3,5) -- (4,5) -- (4,6) -- (3,6) -- cycle}
\fxvcell{3}{6}{4}{7}{0.39}{0}{0}
\fxvliquid{(3,6) -- (4,6) -- (4,6.455724) -- (3,6.320332) -- cycle}
\fxvplic{4}{6.455724}{3}{6.320332}
\fxvcell{4}{5}{4.5}{5.5}{1.00}{1}{0}
\fxvliquid{(4,5) -- (4.5,5) -- (4.5,5.5) -- (4,5.5) -- cycle}
\fxvcell{4.5}{5}{5}{5.5}{1.00}{1}{0}
\fxvliquid{(4.5,5) -- (5,5) -- (5,5.5) -- (4.5,5.5) -- cycle}
\fxvcell{4}{5.5}{4.5}{6}{1.00}{1}{0}
\fxvliquid{(4,5.5) -- (4.5,5.5) -- (4.5,6) -- (4,6) -- cycle}
\fxvcell{4.5}{5.5}{5}{6}{1.00}{1}{0}
\fxvliquid{(4.5,5.5) -- (5,5.5) -- (5,6) -- (4.5,6) -- cycle}
\fxvcell{4}{6}{4.5}{6.5}{1.00}{1}{0}
\fxvliquid{(4,6) -- (4.5,6) -- (4.5,6.5) -- (4,6.5) -- cycle}
\fxvcell{4.5}{6}{5}{6.5}{0.98}{1}{0}
\fxvliquid{(4.5,6) -- (5,6) -- (5,6.451796) -- (4.760683,6.5) -- (4.5,6.5) -- cycle}
\fxvplic{5}{6.451796}{4.760683}{6.5}
\fxvcell{4}{6.5}{4.5}{7}{0.09}{0}{0}
\fxvliquid{(4,6.5) -- (4.5,6.5) -- (4.5,6.565009) -- (4,6.52504) -- cycle}
\fxvplic{4.5}{6.565009}{4}{6.52504}
\fxvcell{4.5}{6.5}{5}{7}{0.03}{0}{0}
\fxvliquid{(4.5,6.5) -- (4.757451,6.5) -- (4.5,6.551857) -- cycle}
\fxvplic{4.757451}{6.5}{4.5}{6.551857}
\fxvdr{(3.854498,5) -- (4,5) -- (4,5.5) -- (3.854498,5.5) -- cycle}{(3.854498,5) -- (4,5) -- (4,5.5) -- (3.854498,5.5) -- cycle}{(3.854498,5) -- (4,5) -- (4,5.5) -- (3.854498,5.5) -- cycle}{0.15}
\fxvdr{(3.863107,5.5) -- (4,5.5) -- (4,6) -- (3.863107,6) -- cycle}{(3.863107,5.5) -- (4,5.5) -- (4,6) -- (3.863107,6) -- cycle}{(3.863107,5.5) -- (4,5.5) -- (4,6) -- (3.863107,6) -- cycle}{0.14}
\fxvdr{(3.863107,6) -- (4,6) -- (4,6.5) -- (3.863107,6.5) -- cycle}{(3.863107,6) -- (4,6) -- (4,6.5) -- (3.863107,6.5) -- cycle}{(3.863107,6) -- (4,6) -- (4,6.455724) -- (3.863107,6.43719) -- cycle}{0.14}
\fxvdr{(3.854498,6.5) -- (4,6.5) -- (4,7) -- (3.854498,7) -- cycle}{(3.854498,6.5) -- (4,6.5) -- (4,7) -- (3.854498,7) -- cycle}{}{0.15}
\fxvdr{(4.354498,5) -- (4.5,5) -- (4.5,5.5) -- (4.354498,5.5) -- cycle}{(4.354498,5) -- (4.5,5) -- (4.5,5.5) -- (4.354498,5.5) -- cycle}{(4.354498,5) -- (4.5,5) -- (4.5,5.5) -- (4.354498,5.5) -- cycle}{0.29}
\fxvdr{(4.363107,5.5) -- (4.5,5.5) -- (4.5,6) -- (4.363107,6) -- cycle}{(4.363107,5.5) -- (4.5,5.5) -- (4.5,6) -- (4.363107,6) -- cycle}{(4.363107,5.5) -- (4.5,5.5) -- (4.5,6) -- (4.363107,6) -- cycle}{0.27}
\fxvdr{(4.363107,6) -- (4.5,6) -- (4.5,6.5) -- (4.363107,6.5) -- cycle}{(4.363107,6) -- (4.5,6) -- (4.5,6.5) -- (4.363107,6.5) -- cycle}{(4.363107,6) -- (4.5,6) -- (4.5,6.5) -- (4.363107,6.5) -- cycle}{0.27}
\fxvdr{(4.354498,6.5) -- (4.5,6.5) -- (4.5,7) -- (4.354498,7) -- cycle}{(4.354498,6.5) -- (4.5,6.5) -- (4.5,7) -- (4.354498,7) -- cycle}{(4.354498,6.5) -- (4.5,6.5) -- (4.5,6.565009) -- (4.354498,6.553378) -- cycle}{0.29}
\fxvdr{(4.854498,5) -- (5,5) -- (5,5.5) -- (4.854498,5.5) -- cycle}{(4.854498,5) -- (5,5) -- (5,5.5) -- (4.854498,5.5) -- cycle}{(4.854498,5) -- (5,5) -- (5,5.5) -- (4.854498,5.5) -- cycle}{0.29}
\fxvdr{(4.863107,5.5) -- (5,5.5) -- (5,6) -- (4.863107,6) -- cycle}{(4.863107,5.5) -- (5,5.5) -- (5,6) -- (4.863107,6) -- cycle}{(4.863107,5.5) -- (5,5.5) -- (5,6) -- (4.863107,6) -- cycle}{0.27}
\fxvdr{(4.863107,6) -- (5,6) -- (5,6.5) -- (4.863107,6.5) -- cycle}{(4.863107,6) -- (5,6) -- (5,6.5) -- (4.863107,6.5) -- cycle}{(4.863107,6) -- (5,6) -- (5,6.451796) -- (4.863107,6.479369) -- cycle}{0.27}
\fxvdr{(4.854498,6.5) -- (5,6.5) -- (5,7) -- (4.854498,7) -- cycle}{(4.854498,6.5) -- (5,6.5) -- (5,7) -- (4.854498,7) -- cycle}{}{0.29}
\fxvvel{1}{3}{5.5}{0.730}
\fxvvel{1}{3}{6.5}{0.730}
\fxvvel{1}{4}{5.25}{0.752}
\fxvvel{1}{4}{5.75}{0.708}
\fxvvel{1}{4}{6.25}{0.708}
\fxvvel{1}{4}{6.75}{0.752}
\fxvvel{1}{4.5}{5.25}{0.752}
\fxvvel{1}{4.5}{5.75}{0.708}
\fxvvel{1}{4.5}{6.25}{0.708}
\fxvvel{1}{4.5}{6.75}{0.752}
\fxvvel{1}{5}{5.25}{0.752}
\fxvvel{1}{5}{5.75}{0.708}
\fxvvel{1}{5}{6.25}{0.708}
\fxvvel{1}{5}{6.75}{0.752}
\fi
\ifnum\fxvpanel=3
\fxvcell{1}{0}{2}{1}{1.00}{1}{1}
\fxvsolid{(1,0) -- (2,0) -- (2,0.45) -- (1,0.25) -- cycle}
\fxvwall{2}{0.45}{1}{0.25}
\fxvliquid{(2,0.45) -- (2,1) -- (1,1) -- (1,0.25) -- cycle}
\fxvcell{2}{0}{3}{1}{0.66}{1}{1}
\fxvsolid{(2,0) -- (3,0) -- (3,0.55) -- (2,0.35) -- cycle}
\fxvwall{3}{0.55}{2}{0.35}
\fxvliquid{(2.7,1) -- (2,1) -- (2,0.35) -- (2.537234,0.457447) -- cycle}
\fxvplic{2.7}{1}{2.537234}{0.457447}
\fxvcell{3}{0}{4}{1}{0.00}{0}{1}
\fxvsolid{(3,0) -- (4,0) -- (4,0.65) -- (3,0.45) -- cycle}
\fxvwall{4}{0.65}{3}{0.45}
\fxvcell{1}{1}{2}{2}{1.00}{1}{0}
\fxvliquid{(1,1) -- (2,1) -- (2,2) -- (1,2) -- cycle}
\fxvcell{2}{1}{3}{2}{0.85}{1}{0}
\fxvliquid{(2,1) -- (2.7,1) -- (3,2) -- (2,2) -- cycle}
\fxvplic{2.7}{1}{3}{2}
\fxvcell{3}{1}{4}{2}{0.00}{0}{0}
\fxvdr{(1.577856,0) -- (2,0) -- (2,1) -- (1.577856,1) -- cycle}{(2,0.45) -- (2,1) -- (1.577856,1) -- (1.577856,0.365571) -- cycle}{(2,0.45) -- (2,1) -- (1.577856,1) -- (1.577856,0.365571) -- cycle}{0.42}
\fxvdr{(1.583333,1) -- (2,1) -- (2,2) -- (1.583333,2) -- cycle}{(1.583333,1) -- (2,1) -- (2,2) -- (1.583333,2) -- cycle}{(1.583333,1) -- (2,1) -- (2,2) -- (1.583333,2) -- cycle}{0.42}
\fxvdr{(2.576831,0) -- (3,0) -- (3,1) -- (2.576831,1) -- cycle}{(3,0.55) -- (3,1) -- (2.576831,1) -- (2.576831,0.465366) -- cycle}{(2.7,1) -- (2.576831,1) -- (2.576831,0.589436) -- cycle}{0.42}
\fxvdr{(2.583333,1) -- (3,1) -- (3,2) -- (2.583333,2) -- cycle}{(2.583333,1) -- (3,1) -- (3,2) -- (2.583333,2) -- cycle}{(2.583333,1) -- (2.7,1) -- (3,2) -- (2.583333,2) -- cycle}{0.42}
\fxvdr{(3.575335,0) -- (4,0) -- (4,1) -- (3.575335,1) -- cycle}{(4,0.65) -- (4,1) -- (3.575335,1) -- (3.575335,0.565067) -- cycle}{}{0.42}
\fxvdr{(3.583333,1) -- (4,1) -- (4,2) -- (3.583333,2) -- cycle}{(3.583333,1) -- (4,1) -- (4,2) -- (3.583333,2) -- cycle}{}{0.42}
\fxvvel{1}{1}{0.5}{1.000}
\fxvvel{1}{1}{1.5}{1.000}
\fxvvel{1}{2}{0.5}{1.000}
\fxvvel{1}{2}{1.5}{1.000}
\fxvvel{1}{3}{0.5}{1.000}
\fxvvel{1}{3}{1.5}{1.000}
\fxvvel{1}{4}{0.5}{1.000}
\fxvvel{1}{4}{1.5}{1.000}
\fi
%
    \fxvnone
    \def\fxvcell##1##2##3##4##5##6##7{\draw[tr cell] (##1,##2) rectangle (##3,##4);}%
    \def\fxvplic##1##2##3##4{\draw[tr plic] (##1,##2) -- (##3,##4);}%
    \def\fxvwall##1##2##3##4{\draw[tr solid line] (##1,##2) -- (##3,##4);}%
    % Generated by docs/techref/figures/generate.jl; do not edit.
\ifnum\fxvpanel=1
\fxvcell{0}{3}{1}{4}{0.00}{0}{0}
\fxvcell{1}{3}{2}{4}{0.01}{0}{0}
\fxvliquid{(2,3.811366) -- (2,4) -- (1.929225,4) -- cycle}
\fxvplic{2}{3.811366}{1.929225}{4}
\fxvcell{2}{3}{3}{4}{0.92}{1}{0}
\fxvliquid{(2.333083,3) -- (3,3) -- (3,4) -- (2,4) -- (2,3.50418) -- cycle}
\fxvplic{2.333083}{3}{2}{3.50418}
\fxvcell{0}{4}{1}{5}{0.00}{0}{0}
\fxvcell{1}{4}{2}{5}{0.03}{0}{0}
\fxvliquid{(1.912944,4) -- (2,4) -- (2,4.58673) -- cycle}
\fxvplic{1.912944}{4}{2}{4.58673}
\fxvcell{2}{4}{3}{5}{0.99}{1}{0}
\fxvliquid{(2,4) -- (3,4) -- (3,5) -- (2.100546,5) -- (2,4.739856) -- cycle}
\fxvplic{2.100546}{5}{2}{4.739856}
\fxvcell{0}{5}{1}{6}{0.00}{0}{0}
\fxvcell{1}{5}{2}{6}{0.00}{0}{0}
\fxvcell{2}{5}{3}{6}{0.63}{1}{0}
\fxvliquid{(2,5) -- (3,5) -- (3,6) -- (2.955594,6) -- (2,5.217152) -- cycle}
\fxvplic{2.955594}{6}{2}{5.217152}
\fxvdr{(0.784912,3) -- (1,3) -- (1,4) -- (0.784912,4) -- cycle}{(0.784912,3) -- (1,3) -- (1,4) -- (0.784912,4) -- cycle}{}{0.22}
\fxvdr{(0.828196,4) -- (1,4) -- (1,5) -- (0.828196,5) -- cycle}{(0.828196,4) -- (1,4) -- (1,5) -- (0.828196,5) -- cycle}{}{0.17}
\fxvdr{(0.858803,5) -- (1,5) -- (1,6) -- (0.858803,6) -- cycle}{(0.858803,5) -- (1,5) -- (1,6) -- (0.858803,6) -- cycle}{}{0.14}
\fxvdr{(1.784912,3) -- (2,3) -- (2,4) -- (1.784912,4) -- cycle}{(1.784912,3) -- (2,3) -- (2,4) -- (1.784912,4) -- cycle}{(2,3.811366) -- (2,4) -- (1.929225,4) -- cycle}{0.22}
\fxvdr{(1.828196,4) -- (2,4) -- (2,5) -- (1.828196,5) -- cycle}{(1.828196,4) -- (2,4) -- (2,5) -- (1.828196,5) -- cycle}{(1.912944,4) -- (2,4) -- (2,4.58673) -- cycle}{0.17}
\fxvdr{(1.858803,5) -- (2,5) -- (2,6) -- (1.858803,6) -- cycle}{(1.858803,5) -- (2,5) -- (2,6) -- (1.858803,6) -- cycle}{}{0.14}
\fxvdr{(2.784912,3) -- (3,3) -- (3,4) -- (2.784912,4) -- cycle}{(2.784912,3) -- (3,3) -- (3,4) -- (2.784912,4) -- cycle}{(2.784912,3) -- (3,3) -- (3,4) -- (2.784912,4) -- cycle}{0.22}
\fxvdr{(2.828196,4) -- (3,4) -- (3,5) -- (2.828196,5) -- cycle}{(2.828196,4) -- (3,4) -- (3,5) -- (2.828196,5) -- cycle}{(2.828196,4) -- (3,4) -- (3,5) -- (2.828196,5) -- cycle}{0.17}
\fxvdr{(2.858803,5) -- (3,5) -- (3,6) -- (2.858803,6) -- cycle}{(2.858803,5) -- (3,5) -- (3,6) -- (2.858803,6) -- cycle}{(2.858803,5) -- (3,5) -- (3,6) -- (2.955594,6) -- (2.858803,5.920706) -- cycle}{0.14}
\fxvvel{1}{0}{3.5}{1.112}
\fxvvel{1}{0}{4.5}{0.888}
\fxvvel{1}{0}{5.5}{0.730}
\fxvvel{1}{1}{3.5}{1.112}
\fxvvel{1}{1}{4.5}{0.888}
\fxvvel{1}{1}{5.5}{0.730}
\fxvvel{1}{2}{3.5}{1.112}
\fxvvel{1}{2}{4.5}{0.888}
\fxvvel{1}{2}{5.5}{0.730}
\fxvvel{1}{3}{3.5}{1.112}
\fxvvel{1}{3}{4.5}{0.888}
\fxvvel{1}{3}{5.5}{0.730}
\fi
\ifnum\fxvpanel=2
\fxvcell{3}{5}{4}{6}{1.00}{1}{0}
\fxvliquid{(3,5) -- (4,5) -- (4,6) -- (3,6) -- cycle}
\fxvcell{3}{6}{4}{7}{0.39}{0}{0}
\fxvliquid{(3,6) -- (4,6) -- (4,6.455724) -- (3,6.320332) -- cycle}
\fxvplic{4}{6.455724}{3}{6.320332}
\fxvcell{4}{5}{4.5}{5.5}{1.00}{1}{0}
\fxvliquid{(4,5) -- (4.5,5) -- (4.5,5.5) -- (4,5.5) -- cycle}
\fxvcell{4.5}{5}{5}{5.5}{1.00}{1}{0}
\fxvliquid{(4.5,5) -- (5,5) -- (5,5.5) -- (4.5,5.5) -- cycle}
\fxvcell{4}{5.5}{4.5}{6}{1.00}{1}{0}
\fxvliquid{(4,5.5) -- (4.5,5.5) -- (4.5,6) -- (4,6) -- cycle}
\fxvcell{4.5}{5.5}{5}{6}{1.00}{1}{0}
\fxvliquid{(4.5,5.5) -- (5,5.5) -- (5,6) -- (4.5,6) -- cycle}
\fxvcell{4}{6}{4.5}{6.5}{1.00}{1}{0}
\fxvliquid{(4,6) -- (4.5,6) -- (4.5,6.5) -- (4,6.5) -- cycle}
\fxvcell{4.5}{6}{5}{6.5}{0.98}{1}{0}
\fxvliquid{(4.5,6) -- (5,6) -- (5,6.451796) -- (4.760683,6.5) -- (4.5,6.5) -- cycle}
\fxvplic{5}{6.451796}{4.760683}{6.5}
\fxvcell{4}{6.5}{4.5}{7}{0.09}{0}{0}
\fxvliquid{(4,6.5) -- (4.5,6.5) -- (4.5,6.565009) -- (4,6.52504) -- cycle}
\fxvplic{4.5}{6.565009}{4}{6.52504}
\fxvcell{4.5}{6.5}{5}{7}{0.03}{0}{0}
\fxvliquid{(4.5,6.5) -- (4.757451,6.5) -- (4.5,6.551857) -- cycle}
\fxvplic{4.757451}{6.5}{4.5}{6.551857}
\fxvdr{(3.854498,5) -- (4,5) -- (4,5.5) -- (3.854498,5.5) -- cycle}{(3.854498,5) -- (4,5) -- (4,5.5) -- (3.854498,5.5) -- cycle}{(3.854498,5) -- (4,5) -- (4,5.5) -- (3.854498,5.5) -- cycle}{0.15}
\fxvdr{(3.863107,5.5) -- (4,5.5) -- (4,6) -- (3.863107,6) -- cycle}{(3.863107,5.5) -- (4,5.5) -- (4,6) -- (3.863107,6) -- cycle}{(3.863107,5.5) -- (4,5.5) -- (4,6) -- (3.863107,6) -- cycle}{0.14}
\fxvdr{(3.863107,6) -- (4,6) -- (4,6.5) -- (3.863107,6.5) -- cycle}{(3.863107,6) -- (4,6) -- (4,6.5) -- (3.863107,6.5) -- cycle}{(3.863107,6) -- (4,6) -- (4,6.455724) -- (3.863107,6.43719) -- cycle}{0.14}
\fxvdr{(3.854498,6.5) -- (4,6.5) -- (4,7) -- (3.854498,7) -- cycle}{(3.854498,6.5) -- (4,6.5) -- (4,7) -- (3.854498,7) -- cycle}{}{0.15}
\fxvdr{(4.354498,5) -- (4.5,5) -- (4.5,5.5) -- (4.354498,5.5) -- cycle}{(4.354498,5) -- (4.5,5) -- (4.5,5.5) -- (4.354498,5.5) -- cycle}{(4.354498,5) -- (4.5,5) -- (4.5,5.5) -- (4.354498,5.5) -- cycle}{0.29}
\fxvdr{(4.363107,5.5) -- (4.5,5.5) -- (4.5,6) -- (4.363107,6) -- cycle}{(4.363107,5.5) -- (4.5,5.5) -- (4.5,6) -- (4.363107,6) -- cycle}{(4.363107,5.5) -- (4.5,5.5) -- (4.5,6) -- (4.363107,6) -- cycle}{0.27}
\fxvdr{(4.363107,6) -- (4.5,6) -- (4.5,6.5) -- (4.363107,6.5) -- cycle}{(4.363107,6) -- (4.5,6) -- (4.5,6.5) -- (4.363107,6.5) -- cycle}{(4.363107,6) -- (4.5,6) -- (4.5,6.5) -- (4.363107,6.5) -- cycle}{0.27}
\fxvdr{(4.354498,6.5) -- (4.5,6.5) -- (4.5,7) -- (4.354498,7) -- cycle}{(4.354498,6.5) -- (4.5,6.5) -- (4.5,7) -- (4.354498,7) -- cycle}{(4.354498,6.5) -- (4.5,6.5) -- (4.5,6.565009) -- (4.354498,6.553378) -- cycle}{0.29}
\fxvdr{(4.854498,5) -- (5,5) -- (5,5.5) -- (4.854498,5.5) -- cycle}{(4.854498,5) -- (5,5) -- (5,5.5) -- (4.854498,5.5) -- cycle}{(4.854498,5) -- (5,5) -- (5,5.5) -- (4.854498,5.5) -- cycle}{0.29}
\fxvdr{(4.863107,5.5) -- (5,5.5) -- (5,6) -- (4.863107,6) -- cycle}{(4.863107,5.5) -- (5,5.5) -- (5,6) -- (4.863107,6) -- cycle}{(4.863107,5.5) -- (5,5.5) -- (5,6) -- (4.863107,6) -- cycle}{0.27}
\fxvdr{(4.863107,6) -- (5,6) -- (5,6.5) -- (4.863107,6.5) -- cycle}{(4.863107,6) -- (5,6) -- (5,6.5) -- (4.863107,6.5) -- cycle}{(4.863107,6) -- (5,6) -- (5,6.451796) -- (4.863107,6.479369) -- cycle}{0.27}
\fxvdr{(4.854498,6.5) -- (5,6.5) -- (5,7) -- (4.854498,7) -- cycle}{(4.854498,6.5) -- (5,6.5) -- (5,7) -- (4.854498,7) -- cycle}{}{0.29}
\fxvvel{1}{3}{5.5}{0.730}
\fxvvel{1}{3}{6.5}{0.730}
\fxvvel{1}{4}{5.25}{0.752}
\fxvvel{1}{4}{5.75}{0.708}
\fxvvel{1}{4}{6.25}{0.708}
\fxvvel{1}{4}{6.75}{0.752}
\fxvvel{1}{4.5}{5.25}{0.752}
\fxvvel{1}{4.5}{5.75}{0.708}
\fxvvel{1}{4.5}{6.25}{0.708}
\fxvvel{1}{4.5}{6.75}{0.752}
\fxvvel{1}{5}{5.25}{0.752}
\fxvvel{1}{5}{5.75}{0.708}
\fxvvel{1}{5}{6.25}{0.708}
\fxvvel{1}{5}{6.75}{0.752}
\fi
\ifnum\fxvpanel=3
\fxvcell{1}{0}{2}{1}{1.00}{1}{1}
\fxvsolid{(1,0) -- (2,0) -- (2,0.45) -- (1,0.25) -- cycle}
\fxvwall{2}{0.45}{1}{0.25}
\fxvliquid{(2,0.45) -- (2,1) -- (1,1) -- (1,0.25) -- cycle}
\fxvcell{2}{0}{3}{1}{0.66}{1}{1}
\fxvsolid{(2,0) -- (3,0) -- (3,0.55) -- (2,0.35) -- cycle}
\fxvwall{3}{0.55}{2}{0.35}
\fxvliquid{(2.7,1) -- (2,1) -- (2,0.35) -- (2.537234,0.457447) -- cycle}
\fxvplic{2.7}{1}{2.537234}{0.457447}
\fxvcell{3}{0}{4}{1}{0.00}{0}{1}
\fxvsolid{(3,0) -- (4,0) -- (4,0.65) -- (3,0.45) -- cycle}
\fxvwall{4}{0.65}{3}{0.45}
\fxvcell{1}{1}{2}{2}{1.00}{1}{0}
\fxvliquid{(1,1) -- (2,1) -- (2,2) -- (1,2) -- cycle}
\fxvcell{2}{1}{3}{2}{0.85}{1}{0}
\fxvliquid{(2,1) -- (2.7,1) -- (3,2) -- (2,2) -- cycle}
\fxvplic{2.7}{1}{3}{2}
\fxvcell{3}{1}{4}{2}{0.00}{0}{0}
\fxvdr{(1.577856,0) -- (2,0) -- (2,1) -- (1.577856,1) -- cycle}{(2,0.45) -- (2,1) -- (1.577856,1) -- (1.577856,0.365571) -- cycle}{(2,0.45) -- (2,1) -- (1.577856,1) -- (1.577856,0.365571) -- cycle}{0.42}
\fxvdr{(1.583333,1) -- (2,1) -- (2,2) -- (1.583333,2) -- cycle}{(1.583333,1) -- (2,1) -- (2,2) -- (1.583333,2) -- cycle}{(1.583333,1) -- (2,1) -- (2,2) -- (1.583333,2) -- cycle}{0.42}
\fxvdr{(2.576831,0) -- (3,0) -- (3,1) -- (2.576831,1) -- cycle}{(3,0.55) -- (3,1) -- (2.576831,1) -- (2.576831,0.465366) -- cycle}{(2.7,1) -- (2.576831,1) -- (2.576831,0.589436) -- cycle}{0.42}
\fxvdr{(2.583333,1) -- (3,1) -- (3,2) -- (2.583333,2) -- cycle}{(2.583333,1) -- (3,1) -- (3,2) -- (2.583333,2) -- cycle}{(2.583333,1) -- (2.7,1) -- (3,2) -- (2.583333,2) -- cycle}{0.42}
\fxvdr{(3.575335,0) -- (4,0) -- (4,1) -- (3.575335,1) -- cycle}{(4,0.65) -- (4,1) -- (3.575335,1) -- (3.575335,0.565067) -- cycle}{}{0.42}
\fxvdr{(3.583333,1) -- (4,1) -- (4,2) -- (3.583333,2) -- cycle}{(3.583333,1) -- (4,1) -- (4,2) -- (3.583333,2) -- cycle}{}{0.42}
\fxvvel{1}{1}{0.5}{1.000}
\fxvvel{1}{1}{1.5}{1.000}
\fxvvel{1}{2}{0.5}{1.000}
\fxvvel{1}{2}{1.5}{1.000}
\fxvvel{1}{3}{0.5}{1.000}
\fxvvel{1}{3}{1.5}{1.000}
\fxvvel{1}{4}{0.5}{1.000}
\fxvvel{1}{4}{1.5}{1.000}
\fi
%
    \fxvnone
    \def\fxvdr##1##2##3##4{\draw[tr dr] ##1;}%
    % Generated by docs/techref/figures/generate.jl; do not edit.
\ifnum\fxvpanel=1
\fxvcell{0}{3}{1}{4}{0.00}{0}{0}
\fxvcell{1}{3}{2}{4}{0.01}{0}{0}
\fxvliquid{(2,3.811366) -- (2,4) -- (1.929225,4) -- cycle}
\fxvplic{2}{3.811366}{1.929225}{4}
\fxvcell{2}{3}{3}{4}{0.92}{1}{0}
\fxvliquid{(2.333083,3) -- (3,3) -- (3,4) -- (2,4) -- (2,3.50418) -- cycle}
\fxvplic{2.333083}{3}{2}{3.50418}
\fxvcell{0}{4}{1}{5}{0.00}{0}{0}
\fxvcell{1}{4}{2}{5}{0.03}{0}{0}
\fxvliquid{(1.912944,4) -- (2,4) -- (2,4.58673) -- cycle}
\fxvplic{1.912944}{4}{2}{4.58673}
\fxvcell{2}{4}{3}{5}{0.99}{1}{0}
\fxvliquid{(2,4) -- (3,4) -- (3,5) -- (2.100546,5) -- (2,4.739856) -- cycle}
\fxvplic{2.100546}{5}{2}{4.739856}
\fxvcell{0}{5}{1}{6}{0.00}{0}{0}
\fxvcell{1}{5}{2}{6}{0.00}{0}{0}
\fxvcell{2}{5}{3}{6}{0.63}{1}{0}
\fxvliquid{(2,5) -- (3,5) -- (3,6) -- (2.955594,6) -- (2,5.217152) -- cycle}
\fxvplic{2.955594}{6}{2}{5.217152}
\fxvdr{(0.784912,3) -- (1,3) -- (1,4) -- (0.784912,4) -- cycle}{(0.784912,3) -- (1,3) -- (1,4) -- (0.784912,4) -- cycle}{}{0.22}
\fxvdr{(0.828196,4) -- (1,4) -- (1,5) -- (0.828196,5) -- cycle}{(0.828196,4) -- (1,4) -- (1,5) -- (0.828196,5) -- cycle}{}{0.17}
\fxvdr{(0.858803,5) -- (1,5) -- (1,6) -- (0.858803,6) -- cycle}{(0.858803,5) -- (1,5) -- (1,6) -- (0.858803,6) -- cycle}{}{0.14}
\fxvdr{(1.784912,3) -- (2,3) -- (2,4) -- (1.784912,4) -- cycle}{(1.784912,3) -- (2,3) -- (2,4) -- (1.784912,4) -- cycle}{(2,3.811366) -- (2,4) -- (1.929225,4) -- cycle}{0.22}
\fxvdr{(1.828196,4) -- (2,4) -- (2,5) -- (1.828196,5) -- cycle}{(1.828196,4) -- (2,4) -- (2,5) -- (1.828196,5) -- cycle}{(1.912944,4) -- (2,4) -- (2,4.58673) -- cycle}{0.17}
\fxvdr{(1.858803,5) -- (2,5) -- (2,6) -- (1.858803,6) -- cycle}{(1.858803,5) -- (2,5) -- (2,6) -- (1.858803,6) -- cycle}{}{0.14}
\fxvdr{(2.784912,3) -- (3,3) -- (3,4) -- (2.784912,4) -- cycle}{(2.784912,3) -- (3,3) -- (3,4) -- (2.784912,4) -- cycle}{(2.784912,3) -- (3,3) -- (3,4) -- (2.784912,4) -- cycle}{0.22}
\fxvdr{(2.828196,4) -- (3,4) -- (3,5) -- (2.828196,5) -- cycle}{(2.828196,4) -- (3,4) -- (3,5) -- (2.828196,5) -- cycle}{(2.828196,4) -- (3,4) -- (3,5) -- (2.828196,5) -- cycle}{0.17}
\fxvdr{(2.858803,5) -- (3,5) -- (3,6) -- (2.858803,6) -- cycle}{(2.858803,5) -- (3,5) -- (3,6) -- (2.858803,6) -- cycle}{(2.858803,5) -- (3,5) -- (3,6) -- (2.955594,6) -- (2.858803,5.920706) -- cycle}{0.14}
\fxvvel{1}{0}{3.5}{1.112}
\fxvvel{1}{0}{4.5}{0.888}
\fxvvel{1}{0}{5.5}{0.730}
\fxvvel{1}{1}{3.5}{1.112}
\fxvvel{1}{1}{4.5}{0.888}
\fxvvel{1}{1}{5.5}{0.730}
\fxvvel{1}{2}{3.5}{1.112}
\fxvvel{1}{2}{4.5}{0.888}
\fxvvel{1}{2}{5.5}{0.730}
\fxvvel{1}{3}{3.5}{1.112}
\fxvvel{1}{3}{4.5}{0.888}
\fxvvel{1}{3}{5.5}{0.730}
\fi
\ifnum\fxvpanel=2
\fxvcell{3}{5}{4}{6}{1.00}{1}{0}
\fxvliquid{(3,5) -- (4,5) -- (4,6) -- (3,6) -- cycle}
\fxvcell{3}{6}{4}{7}{0.39}{0}{0}
\fxvliquid{(3,6) -- (4,6) -- (4,6.455724) -- (3,6.320332) -- cycle}
\fxvplic{4}{6.455724}{3}{6.320332}
\fxvcell{4}{5}{4.5}{5.5}{1.00}{1}{0}
\fxvliquid{(4,5) -- (4.5,5) -- (4.5,5.5) -- (4,5.5) -- cycle}
\fxvcell{4.5}{5}{5}{5.5}{1.00}{1}{0}
\fxvliquid{(4.5,5) -- (5,5) -- (5,5.5) -- (4.5,5.5) -- cycle}
\fxvcell{4}{5.5}{4.5}{6}{1.00}{1}{0}
\fxvliquid{(4,5.5) -- (4.5,5.5) -- (4.5,6) -- (4,6) -- cycle}
\fxvcell{4.5}{5.5}{5}{6}{1.00}{1}{0}
\fxvliquid{(4.5,5.5) -- (5,5.5) -- (5,6) -- (4.5,6) -- cycle}
\fxvcell{4}{6}{4.5}{6.5}{1.00}{1}{0}
\fxvliquid{(4,6) -- (4.5,6) -- (4.5,6.5) -- (4,6.5) -- cycle}
\fxvcell{4.5}{6}{5}{6.5}{0.98}{1}{0}
\fxvliquid{(4.5,6) -- (5,6) -- (5,6.451796) -- (4.760683,6.5) -- (4.5,6.5) -- cycle}
\fxvplic{5}{6.451796}{4.760683}{6.5}
\fxvcell{4}{6.5}{4.5}{7}{0.09}{0}{0}
\fxvliquid{(4,6.5) -- (4.5,6.5) -- (4.5,6.565009) -- (4,6.52504) -- cycle}
\fxvplic{4.5}{6.565009}{4}{6.52504}
\fxvcell{4.5}{6.5}{5}{7}{0.03}{0}{0}
\fxvliquid{(4.5,6.5) -- (4.757451,6.5) -- (4.5,6.551857) -- cycle}
\fxvplic{4.757451}{6.5}{4.5}{6.551857}
\fxvdr{(3.854498,5) -- (4,5) -- (4,5.5) -- (3.854498,5.5) -- cycle}{(3.854498,5) -- (4,5) -- (4,5.5) -- (3.854498,5.5) -- cycle}{(3.854498,5) -- (4,5) -- (4,5.5) -- (3.854498,5.5) -- cycle}{0.15}
\fxvdr{(3.863107,5.5) -- (4,5.5) -- (4,6) -- (3.863107,6) -- cycle}{(3.863107,5.5) -- (4,5.5) -- (4,6) -- (3.863107,6) -- cycle}{(3.863107,5.5) -- (4,5.5) -- (4,6) -- (3.863107,6) -- cycle}{0.14}
\fxvdr{(3.863107,6) -- (4,6) -- (4,6.5) -- (3.863107,6.5) -- cycle}{(3.863107,6) -- (4,6) -- (4,6.5) -- (3.863107,6.5) -- cycle}{(3.863107,6) -- (4,6) -- (4,6.455724) -- (3.863107,6.43719) -- cycle}{0.14}
\fxvdr{(3.854498,6.5) -- (4,6.5) -- (4,7) -- (3.854498,7) -- cycle}{(3.854498,6.5) -- (4,6.5) -- (4,7) -- (3.854498,7) -- cycle}{}{0.15}
\fxvdr{(4.354498,5) -- (4.5,5) -- (4.5,5.5) -- (4.354498,5.5) -- cycle}{(4.354498,5) -- (4.5,5) -- (4.5,5.5) -- (4.354498,5.5) -- cycle}{(4.354498,5) -- (4.5,5) -- (4.5,5.5) -- (4.354498,5.5) -- cycle}{0.29}
\fxvdr{(4.363107,5.5) -- (4.5,5.5) -- (4.5,6) -- (4.363107,6) -- cycle}{(4.363107,5.5) -- (4.5,5.5) -- (4.5,6) -- (4.363107,6) -- cycle}{(4.363107,5.5) -- (4.5,5.5) -- (4.5,6) -- (4.363107,6) -- cycle}{0.27}
\fxvdr{(4.363107,6) -- (4.5,6) -- (4.5,6.5) -- (4.363107,6.5) -- cycle}{(4.363107,6) -- (4.5,6) -- (4.5,6.5) -- (4.363107,6.5) -- cycle}{(4.363107,6) -- (4.5,6) -- (4.5,6.5) -- (4.363107,6.5) -- cycle}{0.27}
\fxvdr{(4.354498,6.5) -- (4.5,6.5) -- (4.5,7) -- (4.354498,7) -- cycle}{(4.354498,6.5) -- (4.5,6.5) -- (4.5,7) -- (4.354498,7) -- cycle}{(4.354498,6.5) -- (4.5,6.5) -- (4.5,6.565009) -- (4.354498,6.553378) -- cycle}{0.29}
\fxvdr{(4.854498,5) -- (5,5) -- (5,5.5) -- (4.854498,5.5) -- cycle}{(4.854498,5) -- (5,5) -- (5,5.5) -- (4.854498,5.5) -- cycle}{(4.854498,5) -- (5,5) -- (5,5.5) -- (4.854498,5.5) -- cycle}{0.29}
\fxvdr{(4.863107,5.5) -- (5,5.5) -- (5,6) -- (4.863107,6) -- cycle}{(4.863107,5.5) -- (5,5.5) -- (5,6) -- (4.863107,6) -- cycle}{(4.863107,5.5) -- (5,5.5) -- (5,6) -- (4.863107,6) -- cycle}{0.27}
\fxvdr{(4.863107,6) -- (5,6) -- (5,6.5) -- (4.863107,6.5) -- cycle}{(4.863107,6) -- (5,6) -- (5,6.5) -- (4.863107,6.5) -- cycle}{(4.863107,6) -- (5,6) -- (5,6.451796) -- (4.863107,6.479369) -- cycle}{0.27}
\fxvdr{(4.854498,6.5) -- (5,6.5) -- (5,7) -- (4.854498,7) -- cycle}{(4.854498,6.5) -- (5,6.5) -- (5,7) -- (4.854498,7) -- cycle}{}{0.29}
\fxvvel{1}{3}{5.5}{0.730}
\fxvvel{1}{3}{6.5}{0.730}
\fxvvel{1}{4}{5.25}{0.752}
\fxvvel{1}{4}{5.75}{0.708}
\fxvvel{1}{4}{6.25}{0.708}
\fxvvel{1}{4}{6.75}{0.752}
\fxvvel{1}{4.5}{5.25}{0.752}
\fxvvel{1}{4.5}{5.75}{0.708}
\fxvvel{1}{4.5}{6.25}{0.708}
\fxvvel{1}{4.5}{6.75}{0.752}
\fxvvel{1}{5}{5.25}{0.752}
\fxvvel{1}{5}{5.75}{0.708}
\fxvvel{1}{5}{6.25}{0.708}
\fxvvel{1}{5}{6.75}{0.752}
\fi
\ifnum\fxvpanel=3
\fxvcell{1}{0}{2}{1}{1.00}{1}{1}
\fxvsolid{(1,0) -- (2,0) -- (2,0.45) -- (1,0.25) -- cycle}
\fxvwall{2}{0.45}{1}{0.25}
\fxvliquid{(2,0.45) -- (2,1) -- (1,1) -- (1,0.25) -- cycle}
\fxvcell{2}{0}{3}{1}{0.66}{1}{1}
\fxvsolid{(2,0) -- (3,0) -- (3,0.55) -- (2,0.35) -- cycle}
\fxvwall{3}{0.55}{2}{0.35}
\fxvliquid{(2.7,1) -- (2,1) -- (2,0.35) -- (2.537234,0.457447) -- cycle}
\fxvplic{2.7}{1}{2.537234}{0.457447}
\fxvcell{3}{0}{4}{1}{0.00}{0}{1}
\fxvsolid{(3,0) -- (4,0) -- (4,0.65) -- (3,0.45) -- cycle}
\fxvwall{4}{0.65}{3}{0.45}
\fxvcell{1}{1}{2}{2}{1.00}{1}{0}
\fxvliquid{(1,1) -- (2,1) -- (2,2) -- (1,2) -- cycle}
\fxvcell{2}{1}{3}{2}{0.85}{1}{0}
\fxvliquid{(2,1) -- (2.7,1) -- (3,2) -- (2,2) -- cycle}
\fxvplic{2.7}{1}{3}{2}
\fxvcell{3}{1}{4}{2}{0.00}{0}{0}
\fxvdr{(1.577856,0) -- (2,0) -- (2,1) -- (1.577856,1) -- cycle}{(2,0.45) -- (2,1) -- (1.577856,1) -- (1.577856,0.365571) -- cycle}{(2,0.45) -- (2,1) -- (1.577856,1) -- (1.577856,0.365571) -- cycle}{0.42}
\fxvdr{(1.583333,1) -- (2,1) -- (2,2) -- (1.583333,2) -- cycle}{(1.583333,1) -- (2,1) -- (2,2) -- (1.583333,2) -- cycle}{(1.583333,1) -- (2,1) -- (2,2) -- (1.583333,2) -- cycle}{0.42}
\fxvdr{(2.576831,0) -- (3,0) -- (3,1) -- (2.576831,1) -- cycle}{(3,0.55) -- (3,1) -- (2.576831,1) -- (2.576831,0.465366) -- cycle}{(2.7,1) -- (2.576831,1) -- (2.576831,0.589436) -- cycle}{0.42}
\fxvdr{(2.583333,1) -- (3,1) -- (3,2) -- (2.583333,2) -- cycle}{(2.583333,1) -- (3,1) -- (3,2) -- (2.583333,2) -- cycle}{(2.583333,1) -- (2.7,1) -- (3,2) -- (2.583333,2) -- cycle}{0.42}
\fxvdr{(3.575335,0) -- (4,0) -- (4,1) -- (3.575335,1) -- cycle}{(4,0.65) -- (4,1) -- (3.575335,1) -- (3.575335,0.565067) -- cycle}{}{0.42}
\fxvdr{(3.583333,1) -- (4,1) -- (4,2) -- (3.583333,2) -- cycle}{(3.583333,1) -- (4,1) -- (4,2) -- (3.583333,2) -- cycle}{}{0.42}
\fxvvel{1}{1}{0.5}{1.000}
\fxvvel{1}{1}{1.5}{1.000}
\fxvvel{1}{2}{0.5}{1.000}
\fxvvel{1}{2}{1.5}{1.000}
\fxvvel{1}{3}{0.5}{1.000}
\fxvvel{1}{3}{1.5}{1.000}
\fxvvel{1}{4}{0.5}{1.000}
\fxvvel{1}{4}{1.5}{1.000}
\fi
%
    \fxvnone
    \def\fxvcell##1##2##3##4##5##6##7{\ifnum##6=1
      \node[tr chi] at ($(##1,##4)+(0.09,-0.09)$) {};\fi}%
    % Generated by docs/techref/figures/generate.jl; do not edit.
\ifnum\fxvpanel=1
\fxvcell{0}{3}{1}{4}{0.00}{0}{0}
\fxvcell{1}{3}{2}{4}{0.01}{0}{0}
\fxvliquid{(2,3.811366) -- (2,4) -- (1.929225,4) -- cycle}
\fxvplic{2}{3.811366}{1.929225}{4}
\fxvcell{2}{3}{3}{4}{0.92}{1}{0}
\fxvliquid{(2.333083,3) -- (3,3) -- (3,4) -- (2,4) -- (2,3.50418) -- cycle}
\fxvplic{2.333083}{3}{2}{3.50418}
\fxvcell{0}{4}{1}{5}{0.00}{0}{0}
\fxvcell{1}{4}{2}{5}{0.03}{0}{0}
\fxvliquid{(1.912944,4) -- (2,4) -- (2,4.58673) -- cycle}
\fxvplic{1.912944}{4}{2}{4.58673}
\fxvcell{2}{4}{3}{5}{0.99}{1}{0}
\fxvliquid{(2,4) -- (3,4) -- (3,5) -- (2.100546,5) -- (2,4.739856) -- cycle}
\fxvplic{2.100546}{5}{2}{4.739856}
\fxvcell{0}{5}{1}{6}{0.00}{0}{0}
\fxvcell{1}{5}{2}{6}{0.00}{0}{0}
\fxvcell{2}{5}{3}{6}{0.63}{1}{0}
\fxvliquid{(2,5) -- (3,5) -- (3,6) -- (2.955594,6) -- (2,5.217152) -- cycle}
\fxvplic{2.955594}{6}{2}{5.217152}
\fxvdr{(0.784912,3) -- (1,3) -- (1,4) -- (0.784912,4) -- cycle}{(0.784912,3) -- (1,3) -- (1,4) -- (0.784912,4) -- cycle}{}{0.22}
\fxvdr{(0.828196,4) -- (1,4) -- (1,5) -- (0.828196,5) -- cycle}{(0.828196,4) -- (1,4) -- (1,5) -- (0.828196,5) -- cycle}{}{0.17}
\fxvdr{(0.858803,5) -- (1,5) -- (1,6) -- (0.858803,6) -- cycle}{(0.858803,5) -- (1,5) -- (1,6) -- (0.858803,6) -- cycle}{}{0.14}
\fxvdr{(1.784912,3) -- (2,3) -- (2,4) -- (1.784912,4) -- cycle}{(1.784912,3) -- (2,3) -- (2,4) -- (1.784912,4) -- cycle}{(2,3.811366) -- (2,4) -- (1.929225,4) -- cycle}{0.22}
\fxvdr{(1.828196,4) -- (2,4) -- (2,5) -- (1.828196,5) -- cycle}{(1.828196,4) -- (2,4) -- (2,5) -- (1.828196,5) -- cycle}{(1.912944,4) -- (2,4) -- (2,4.58673) -- cycle}{0.17}
\fxvdr{(1.858803,5) -- (2,5) -- (2,6) -- (1.858803,6) -- cycle}{(1.858803,5) -- (2,5) -- (2,6) -- (1.858803,6) -- cycle}{}{0.14}
\fxvdr{(2.784912,3) -- (3,3) -- (3,4) -- (2.784912,4) -- cycle}{(2.784912,3) -- (3,3) -- (3,4) -- (2.784912,4) -- cycle}{(2.784912,3) -- (3,3) -- (3,4) -- (2.784912,4) -- cycle}{0.22}
\fxvdr{(2.828196,4) -- (3,4) -- (3,5) -- (2.828196,5) -- cycle}{(2.828196,4) -- (3,4) -- (3,5) -- (2.828196,5) -- cycle}{(2.828196,4) -- (3,4) -- (3,5) -- (2.828196,5) -- cycle}{0.17}
\fxvdr{(2.858803,5) -- (3,5) -- (3,6) -- (2.858803,6) -- cycle}{(2.858803,5) -- (3,5) -- (3,6) -- (2.858803,6) -- cycle}{(2.858803,5) -- (3,5) -- (3,6) -- (2.955594,6) -- (2.858803,5.920706) -- cycle}{0.14}
\fxvvel{1}{0}{3.5}{1.112}
\fxvvel{1}{0}{4.5}{0.888}
\fxvvel{1}{0}{5.5}{0.730}
\fxvvel{1}{1}{3.5}{1.112}
\fxvvel{1}{1}{4.5}{0.888}
\fxvvel{1}{1}{5.5}{0.730}
\fxvvel{1}{2}{3.5}{1.112}
\fxvvel{1}{2}{4.5}{0.888}
\fxvvel{1}{2}{5.5}{0.730}
\fxvvel{1}{3}{3.5}{1.112}
\fxvvel{1}{3}{4.5}{0.888}
\fxvvel{1}{3}{5.5}{0.730}
\fi
\ifnum\fxvpanel=2
\fxvcell{3}{5}{4}{6}{1.00}{1}{0}
\fxvliquid{(3,5) -- (4,5) -- (4,6) -- (3,6) -- cycle}
\fxvcell{3}{6}{4}{7}{0.39}{0}{0}
\fxvliquid{(3,6) -- (4,6) -- (4,6.455724) -- (3,6.320332) -- cycle}
\fxvplic{4}{6.455724}{3}{6.320332}
\fxvcell{4}{5}{4.5}{5.5}{1.00}{1}{0}
\fxvliquid{(4,5) -- (4.5,5) -- (4.5,5.5) -- (4,5.5) -- cycle}
\fxvcell{4.5}{5}{5}{5.5}{1.00}{1}{0}
\fxvliquid{(4.5,5) -- (5,5) -- (5,5.5) -- (4.5,5.5) -- cycle}
\fxvcell{4}{5.5}{4.5}{6}{1.00}{1}{0}
\fxvliquid{(4,5.5) -- (4.5,5.5) -- (4.5,6) -- (4,6) -- cycle}
\fxvcell{4.5}{5.5}{5}{6}{1.00}{1}{0}
\fxvliquid{(4.5,5.5) -- (5,5.5) -- (5,6) -- (4.5,6) -- cycle}
\fxvcell{4}{6}{4.5}{6.5}{1.00}{1}{0}
\fxvliquid{(4,6) -- (4.5,6) -- (4.5,6.5) -- (4,6.5) -- cycle}
\fxvcell{4.5}{6}{5}{6.5}{0.98}{1}{0}
\fxvliquid{(4.5,6) -- (5,6) -- (5,6.451796) -- (4.760683,6.5) -- (4.5,6.5) -- cycle}
\fxvplic{5}{6.451796}{4.760683}{6.5}
\fxvcell{4}{6.5}{4.5}{7}{0.09}{0}{0}
\fxvliquid{(4,6.5) -- (4.5,6.5) -- (4.5,6.565009) -- (4,6.52504) -- cycle}
\fxvplic{4.5}{6.565009}{4}{6.52504}
\fxvcell{4.5}{6.5}{5}{7}{0.03}{0}{0}
\fxvliquid{(4.5,6.5) -- (4.757451,6.5) -- (4.5,6.551857) -- cycle}
\fxvplic{4.757451}{6.5}{4.5}{6.551857}
\fxvdr{(3.854498,5) -- (4,5) -- (4,5.5) -- (3.854498,5.5) -- cycle}{(3.854498,5) -- (4,5) -- (4,5.5) -- (3.854498,5.5) -- cycle}{(3.854498,5) -- (4,5) -- (4,5.5) -- (3.854498,5.5) -- cycle}{0.15}
\fxvdr{(3.863107,5.5) -- (4,5.5) -- (4,6) -- (3.863107,6) -- cycle}{(3.863107,5.5) -- (4,5.5) -- (4,6) -- (3.863107,6) -- cycle}{(3.863107,5.5) -- (4,5.5) -- (4,6) -- (3.863107,6) -- cycle}{0.14}
\fxvdr{(3.863107,6) -- (4,6) -- (4,6.5) -- (3.863107,6.5) -- cycle}{(3.863107,6) -- (4,6) -- (4,6.5) -- (3.863107,6.5) -- cycle}{(3.863107,6) -- (4,6) -- (4,6.455724) -- (3.863107,6.43719) -- cycle}{0.14}
\fxvdr{(3.854498,6.5) -- (4,6.5) -- (4,7) -- (3.854498,7) -- cycle}{(3.854498,6.5) -- (4,6.5) -- (4,7) -- (3.854498,7) -- cycle}{}{0.15}
\fxvdr{(4.354498,5) -- (4.5,5) -- (4.5,5.5) -- (4.354498,5.5) -- cycle}{(4.354498,5) -- (4.5,5) -- (4.5,5.5) -- (4.354498,5.5) -- cycle}{(4.354498,5) -- (4.5,5) -- (4.5,5.5) -- (4.354498,5.5) -- cycle}{0.29}
\fxvdr{(4.363107,5.5) -- (4.5,5.5) -- (4.5,6) -- (4.363107,6) -- cycle}{(4.363107,5.5) -- (4.5,5.5) -- (4.5,6) -- (4.363107,6) -- cycle}{(4.363107,5.5) -- (4.5,5.5) -- (4.5,6) -- (4.363107,6) -- cycle}{0.27}
\fxvdr{(4.363107,6) -- (4.5,6) -- (4.5,6.5) -- (4.363107,6.5) -- cycle}{(4.363107,6) -- (4.5,6) -- (4.5,6.5) -- (4.363107,6.5) -- cycle}{(4.363107,6) -- (4.5,6) -- (4.5,6.5) -- (4.363107,6.5) -- cycle}{0.27}
\fxvdr{(4.354498,6.5) -- (4.5,6.5) -- (4.5,7) -- (4.354498,7) -- cycle}{(4.354498,6.5) -- (4.5,6.5) -- (4.5,7) -- (4.354498,7) -- cycle}{(4.354498,6.5) -- (4.5,6.5) -- (4.5,6.565009) -- (4.354498,6.553378) -- cycle}{0.29}
\fxvdr{(4.854498,5) -- (5,5) -- (5,5.5) -- (4.854498,5.5) -- cycle}{(4.854498,5) -- (5,5) -- (5,5.5) -- (4.854498,5.5) -- cycle}{(4.854498,5) -- (5,5) -- (5,5.5) -- (4.854498,5.5) -- cycle}{0.29}
\fxvdr{(4.863107,5.5) -- (5,5.5) -- (5,6) -- (4.863107,6) -- cycle}{(4.863107,5.5) -- (5,5.5) -- (5,6) -- (4.863107,6) -- cycle}{(4.863107,5.5) -- (5,5.5) -- (5,6) -- (4.863107,6) -- cycle}{0.27}
\fxvdr{(4.863107,6) -- (5,6) -- (5,6.5) -- (4.863107,6.5) -- cycle}{(4.863107,6) -- (5,6) -- (5,6.5) -- (4.863107,6.5) -- cycle}{(4.863107,6) -- (5,6) -- (5,6.451796) -- (4.863107,6.479369) -- cycle}{0.27}
\fxvdr{(4.854498,6.5) -- (5,6.5) -- (5,7) -- (4.854498,7) -- cycle}{(4.854498,6.5) -- (5,6.5) -- (5,7) -- (4.854498,7) -- cycle}{}{0.29}
\fxvvel{1}{3}{5.5}{0.730}
\fxvvel{1}{3}{6.5}{0.730}
\fxvvel{1}{4}{5.25}{0.752}
\fxvvel{1}{4}{5.75}{0.708}
\fxvvel{1}{4}{6.25}{0.708}
\fxvvel{1}{4}{6.75}{0.752}
\fxvvel{1}{4.5}{5.25}{0.752}
\fxvvel{1}{4.5}{5.75}{0.708}
\fxvvel{1}{4.5}{6.25}{0.708}
\fxvvel{1}{4.5}{6.75}{0.752}
\fxvvel{1}{5}{5.25}{0.752}
\fxvvel{1}{5}{5.75}{0.708}
\fxvvel{1}{5}{6.25}{0.708}
\fxvvel{1}{5}{6.75}{0.752}
\fi
\ifnum\fxvpanel=3
\fxvcell{1}{0}{2}{1}{1.00}{1}{1}
\fxvsolid{(1,0) -- (2,0) -- (2,0.45) -- (1,0.25) -- cycle}
\fxvwall{2}{0.45}{1}{0.25}
\fxvliquid{(2,0.45) -- (2,1) -- (1,1) -- (1,0.25) -- cycle}
\fxvcell{2}{0}{3}{1}{0.66}{1}{1}
\fxvsolid{(2,0) -- (3,0) -- (3,0.55) -- (2,0.35) -- cycle}
\fxvwall{3}{0.55}{2}{0.35}
\fxvliquid{(2.7,1) -- (2,1) -- (2,0.35) -- (2.537234,0.457447) -- cycle}
\fxvplic{2.7}{1}{2.537234}{0.457447}
\fxvcell{3}{0}{4}{1}{0.00}{0}{1}
\fxvsolid{(3,0) -- (4,0) -- (4,0.65) -- (3,0.45) -- cycle}
\fxvwall{4}{0.65}{3}{0.45}
\fxvcell{1}{1}{2}{2}{1.00}{1}{0}
\fxvliquid{(1,1) -- (2,1) -- (2,2) -- (1,2) -- cycle}
\fxvcell{2}{1}{3}{2}{0.85}{1}{0}
\fxvliquid{(2,1) -- (2.7,1) -- (3,2) -- (2,2) -- cycle}
\fxvplic{2.7}{1}{3}{2}
\fxvcell{3}{1}{4}{2}{0.00}{0}{0}
\fxvdr{(1.577856,0) -- (2,0) -- (2,1) -- (1.577856,1) -- cycle}{(2,0.45) -- (2,1) -- (1.577856,1) -- (1.577856,0.365571) -- cycle}{(2,0.45) -- (2,1) -- (1.577856,1) -- (1.577856,0.365571) -- cycle}{0.42}
\fxvdr{(1.583333,1) -- (2,1) -- (2,2) -- (1.583333,2) -- cycle}{(1.583333,1) -- (2,1) -- (2,2) -- (1.583333,2) -- cycle}{(1.583333,1) -- (2,1) -- (2,2) -- (1.583333,2) -- cycle}{0.42}
\fxvdr{(2.576831,0) -- (3,0) -- (3,1) -- (2.576831,1) -- cycle}{(3,0.55) -- (3,1) -- (2.576831,1) -- (2.576831,0.465366) -- cycle}{(2.7,1) -- (2.576831,1) -- (2.576831,0.589436) -- cycle}{0.42}
\fxvdr{(2.583333,1) -- (3,1) -- (3,2) -- (2.583333,2) -- cycle}{(2.583333,1) -- (3,1) -- (3,2) -- (2.583333,2) -- cycle}{(2.583333,1) -- (2.7,1) -- (3,2) -- (2.583333,2) -- cycle}{0.42}
\fxvdr{(3.575335,0) -- (4,0) -- (4,1) -- (3.575335,1) -- cycle}{(4,0.65) -- (4,1) -- (3.575335,1) -- (3.575335,0.565067) -- cycle}{}{0.42}
\fxvdr{(3.583333,1) -- (4,1) -- (4,2) -- (3.583333,2) -- cycle}{(3.583333,1) -- (4,1) -- (4,2) -- (3.583333,2) -- cycle}{}{0.42}
\fxvvel{1}{1}{0.5}{1.000}
\fxvvel{1}{1}{1.5}{1.000}
\fxvvel{1}{2}{0.5}{1.000}
\fxvvel{1}{2}{1.5}{1.000}
\fxvvel{1}{3}{0.5}{1.000}
\fxvvel{1}{3}{1.5}{1.000}
\fxvvel{1}{4}{0.5}{1.000}
\fxvvel{1}{4}{1.5}{1.000}
\fi
}%
  \def\fxvarrows#1#2{% panel, arrow length per unit velocity
    \def\fxvpanel{#1}\fxvnone
    \def\fxvvel##1##2##3##4{\draw[tr thin vector,trred] (##2,##3) -- ++(#2*##4,0);}%
    % Generated by docs/techref/figures/generate.jl; do not edit.
\ifnum\fxvpanel=1
\fxvcell{0}{3}{1}{4}{0.00}{0}{0}
\fxvcell{1}{3}{2}{4}{0.01}{0}{0}
\fxvliquid{(2,3.811366) -- (2,4) -- (1.929225,4) -- cycle}
\fxvplic{2}{3.811366}{1.929225}{4}
\fxvcell{2}{3}{3}{4}{0.92}{1}{0}
\fxvliquid{(2.333083,3) -- (3,3) -- (3,4) -- (2,4) -- (2,3.50418) -- cycle}
\fxvplic{2.333083}{3}{2}{3.50418}
\fxvcell{0}{4}{1}{5}{0.00}{0}{0}
\fxvcell{1}{4}{2}{5}{0.03}{0}{0}
\fxvliquid{(1.912944,4) -- (2,4) -- (2,4.58673) -- cycle}
\fxvplic{1.912944}{4}{2}{4.58673}
\fxvcell{2}{4}{3}{5}{0.99}{1}{0}
\fxvliquid{(2,4) -- (3,4) -- (3,5) -- (2.100546,5) -- (2,4.739856) -- cycle}
\fxvplic{2.100546}{5}{2}{4.739856}
\fxvcell{0}{5}{1}{6}{0.00}{0}{0}
\fxvcell{1}{5}{2}{6}{0.00}{0}{0}
\fxvcell{2}{5}{3}{6}{0.63}{1}{0}
\fxvliquid{(2,5) -- (3,5) -- (3,6) -- (2.955594,6) -- (2,5.217152) -- cycle}
\fxvplic{2.955594}{6}{2}{5.217152}
\fxvdr{(0.784912,3) -- (1,3) -- (1,4) -- (0.784912,4) -- cycle}{(0.784912,3) -- (1,3) -- (1,4) -- (0.784912,4) -- cycle}{}{0.22}
\fxvdr{(0.828196,4) -- (1,4) -- (1,5) -- (0.828196,5) -- cycle}{(0.828196,4) -- (1,4) -- (1,5) -- (0.828196,5) -- cycle}{}{0.17}
\fxvdr{(0.858803,5) -- (1,5) -- (1,6) -- (0.858803,6) -- cycle}{(0.858803,5) -- (1,5) -- (1,6) -- (0.858803,6) -- cycle}{}{0.14}
\fxvdr{(1.784912,3) -- (2,3) -- (2,4) -- (1.784912,4) -- cycle}{(1.784912,3) -- (2,3) -- (2,4) -- (1.784912,4) -- cycle}{(2,3.811366) -- (2,4) -- (1.929225,4) -- cycle}{0.22}
\fxvdr{(1.828196,4) -- (2,4) -- (2,5) -- (1.828196,5) -- cycle}{(1.828196,4) -- (2,4) -- (2,5) -- (1.828196,5) -- cycle}{(1.912944,4) -- (2,4) -- (2,4.58673) -- cycle}{0.17}
\fxvdr{(1.858803,5) -- (2,5) -- (2,6) -- (1.858803,6) -- cycle}{(1.858803,5) -- (2,5) -- (2,6) -- (1.858803,6) -- cycle}{}{0.14}
\fxvdr{(2.784912,3) -- (3,3) -- (3,4) -- (2.784912,4) -- cycle}{(2.784912,3) -- (3,3) -- (3,4) -- (2.784912,4) -- cycle}{(2.784912,3) -- (3,3) -- (3,4) -- (2.784912,4) -- cycle}{0.22}
\fxvdr{(2.828196,4) -- (3,4) -- (3,5) -- (2.828196,5) -- cycle}{(2.828196,4) -- (3,4) -- (3,5) -- (2.828196,5) -- cycle}{(2.828196,4) -- (3,4) -- (3,5) -- (2.828196,5) -- cycle}{0.17}
\fxvdr{(2.858803,5) -- (3,5) -- (3,6) -- (2.858803,6) -- cycle}{(2.858803,5) -- (3,5) -- (3,6) -- (2.858803,6) -- cycle}{(2.858803,5) -- (3,5) -- (3,6) -- (2.955594,6) -- (2.858803,5.920706) -- cycle}{0.14}
\fxvvel{1}{0}{3.5}{1.112}
\fxvvel{1}{0}{4.5}{0.888}
\fxvvel{1}{0}{5.5}{0.730}
\fxvvel{1}{1}{3.5}{1.112}
\fxvvel{1}{1}{4.5}{0.888}
\fxvvel{1}{1}{5.5}{0.730}
\fxvvel{1}{2}{3.5}{1.112}
\fxvvel{1}{2}{4.5}{0.888}
\fxvvel{1}{2}{5.5}{0.730}
\fxvvel{1}{3}{3.5}{1.112}
\fxvvel{1}{3}{4.5}{0.888}
\fxvvel{1}{3}{5.5}{0.730}
\fi
\ifnum\fxvpanel=2
\fxvcell{3}{5}{4}{6}{1.00}{1}{0}
\fxvliquid{(3,5) -- (4,5) -- (4,6) -- (3,6) -- cycle}
\fxvcell{3}{6}{4}{7}{0.39}{0}{0}
\fxvliquid{(3,6) -- (4,6) -- (4,6.455724) -- (3,6.320332) -- cycle}
\fxvplic{4}{6.455724}{3}{6.320332}
\fxvcell{4}{5}{4.5}{5.5}{1.00}{1}{0}
\fxvliquid{(4,5) -- (4.5,5) -- (4.5,5.5) -- (4,5.5) -- cycle}
\fxvcell{4.5}{5}{5}{5.5}{1.00}{1}{0}
\fxvliquid{(4.5,5) -- (5,5) -- (5,5.5) -- (4.5,5.5) -- cycle}
\fxvcell{4}{5.5}{4.5}{6}{1.00}{1}{0}
\fxvliquid{(4,5.5) -- (4.5,5.5) -- (4.5,6) -- (4,6) -- cycle}
\fxvcell{4.5}{5.5}{5}{6}{1.00}{1}{0}
\fxvliquid{(4.5,5.5) -- (5,5.5) -- (5,6) -- (4.5,6) -- cycle}
\fxvcell{4}{6}{4.5}{6.5}{1.00}{1}{0}
\fxvliquid{(4,6) -- (4.5,6) -- (4.5,6.5) -- (4,6.5) -- cycle}
\fxvcell{4.5}{6}{5}{6.5}{0.98}{1}{0}
\fxvliquid{(4.5,6) -- (5,6) -- (5,6.451796) -- (4.760683,6.5) -- (4.5,6.5) -- cycle}
\fxvplic{5}{6.451796}{4.760683}{6.5}
\fxvcell{4}{6.5}{4.5}{7}{0.09}{0}{0}
\fxvliquid{(4,6.5) -- (4.5,6.5) -- (4.5,6.565009) -- (4,6.52504) -- cycle}
\fxvplic{4.5}{6.565009}{4}{6.52504}
\fxvcell{4.5}{6.5}{5}{7}{0.03}{0}{0}
\fxvliquid{(4.5,6.5) -- (4.757451,6.5) -- (4.5,6.551857) -- cycle}
\fxvplic{4.757451}{6.5}{4.5}{6.551857}
\fxvdr{(3.854498,5) -- (4,5) -- (4,5.5) -- (3.854498,5.5) -- cycle}{(3.854498,5) -- (4,5) -- (4,5.5) -- (3.854498,5.5) -- cycle}{(3.854498,5) -- (4,5) -- (4,5.5) -- (3.854498,5.5) -- cycle}{0.15}
\fxvdr{(3.863107,5.5) -- (4,5.5) -- (4,6) -- (3.863107,6) -- cycle}{(3.863107,5.5) -- (4,5.5) -- (4,6) -- (3.863107,6) -- cycle}{(3.863107,5.5) -- (4,5.5) -- (4,6) -- (3.863107,6) -- cycle}{0.14}
\fxvdr{(3.863107,6) -- (4,6) -- (4,6.5) -- (3.863107,6.5) -- cycle}{(3.863107,6) -- (4,6) -- (4,6.5) -- (3.863107,6.5) -- cycle}{(3.863107,6) -- (4,6) -- (4,6.455724) -- (3.863107,6.43719) -- cycle}{0.14}
\fxvdr{(3.854498,6.5) -- (4,6.5) -- (4,7) -- (3.854498,7) -- cycle}{(3.854498,6.5) -- (4,6.5) -- (4,7) -- (3.854498,7) -- cycle}{}{0.15}
\fxvdr{(4.354498,5) -- (4.5,5) -- (4.5,5.5) -- (4.354498,5.5) -- cycle}{(4.354498,5) -- (4.5,5) -- (4.5,5.5) -- (4.354498,5.5) -- cycle}{(4.354498,5) -- (4.5,5) -- (4.5,5.5) -- (4.354498,5.5) -- cycle}{0.29}
\fxvdr{(4.363107,5.5) -- (4.5,5.5) -- (4.5,6) -- (4.363107,6) -- cycle}{(4.363107,5.5) -- (4.5,5.5) -- (4.5,6) -- (4.363107,6) -- cycle}{(4.363107,5.5) -- (4.5,5.5) -- (4.5,6) -- (4.363107,6) -- cycle}{0.27}
\fxvdr{(4.363107,6) -- (4.5,6) -- (4.5,6.5) -- (4.363107,6.5) -- cycle}{(4.363107,6) -- (4.5,6) -- (4.5,6.5) -- (4.363107,6.5) -- cycle}{(4.363107,6) -- (4.5,6) -- (4.5,6.5) -- (4.363107,6.5) -- cycle}{0.27}
\fxvdr{(4.354498,6.5) -- (4.5,6.5) -- (4.5,7) -- (4.354498,7) -- cycle}{(4.354498,6.5) -- (4.5,6.5) -- (4.5,7) -- (4.354498,7) -- cycle}{(4.354498,6.5) -- (4.5,6.5) -- (4.5,6.565009) -- (4.354498,6.553378) -- cycle}{0.29}
\fxvdr{(4.854498,5) -- (5,5) -- (5,5.5) -- (4.854498,5.5) -- cycle}{(4.854498,5) -- (5,5) -- (5,5.5) -- (4.854498,5.5) -- cycle}{(4.854498,5) -- (5,5) -- (5,5.5) -- (4.854498,5.5) -- cycle}{0.29}
\fxvdr{(4.863107,5.5) -- (5,5.5) -- (5,6) -- (4.863107,6) -- cycle}{(4.863107,5.5) -- (5,5.5) -- (5,6) -- (4.863107,6) -- cycle}{(4.863107,5.5) -- (5,5.5) -- (5,6) -- (4.863107,6) -- cycle}{0.27}
\fxvdr{(4.863107,6) -- (5,6) -- (5,6.5) -- (4.863107,6.5) -- cycle}{(4.863107,6) -- (5,6) -- (5,6.5) -- (4.863107,6.5) -- cycle}{(4.863107,6) -- (5,6) -- (5,6.451796) -- (4.863107,6.479369) -- cycle}{0.27}
\fxvdr{(4.854498,6.5) -- (5,6.5) -- (5,7) -- (4.854498,7) -- cycle}{(4.854498,6.5) -- (5,6.5) -- (5,7) -- (4.854498,7) -- cycle}{}{0.29}
\fxvvel{1}{3}{5.5}{0.730}
\fxvvel{1}{3}{6.5}{0.730}
\fxvvel{1}{4}{5.25}{0.752}
\fxvvel{1}{4}{5.75}{0.708}
\fxvvel{1}{4}{6.25}{0.708}
\fxvvel{1}{4}{6.75}{0.752}
\fxvvel{1}{4.5}{5.25}{0.752}
\fxvvel{1}{4.5}{5.75}{0.708}
\fxvvel{1}{4.5}{6.25}{0.708}
\fxvvel{1}{4.5}{6.75}{0.752}
\fxvvel{1}{5}{5.25}{0.752}
\fxvvel{1}{5}{5.75}{0.708}
\fxvvel{1}{5}{6.25}{0.708}
\fxvvel{1}{5}{6.75}{0.752}
\fi
\ifnum\fxvpanel=3
\fxvcell{1}{0}{2}{1}{1.00}{1}{1}
\fxvsolid{(1,0) -- (2,0) -- (2,0.45) -- (1,0.25) -- cycle}
\fxvwall{2}{0.45}{1}{0.25}
\fxvliquid{(2,0.45) -- (2,1) -- (1,1) -- (1,0.25) -- cycle}
\fxvcell{2}{0}{3}{1}{0.66}{1}{1}
\fxvsolid{(2,0) -- (3,0) -- (3,0.55) -- (2,0.35) -- cycle}
\fxvwall{3}{0.55}{2}{0.35}
\fxvliquid{(2.7,1) -- (2,1) -- (2,0.35) -- (2.537234,0.457447) -- cycle}
\fxvplic{2.7}{1}{2.537234}{0.457447}
\fxvcell{3}{0}{4}{1}{0.00}{0}{1}
\fxvsolid{(3,0) -- (4,0) -- (4,0.65) -- (3,0.45) -- cycle}
\fxvwall{4}{0.65}{3}{0.45}
\fxvcell{1}{1}{2}{2}{1.00}{1}{0}
\fxvliquid{(1,1) -- (2,1) -- (2,2) -- (1,2) -- cycle}
\fxvcell{2}{1}{3}{2}{0.85}{1}{0}
\fxvliquid{(2,1) -- (2.7,1) -- (3,2) -- (2,2) -- cycle}
\fxvplic{2.7}{1}{3}{2}
\fxvcell{3}{1}{4}{2}{0.00}{0}{0}
\fxvdr{(1.577856,0) -- (2,0) -- (2,1) -- (1.577856,1) -- cycle}{(2,0.45) -- (2,1) -- (1.577856,1) -- (1.577856,0.365571) -- cycle}{(2,0.45) -- (2,1) -- (1.577856,1) -- (1.577856,0.365571) -- cycle}{0.42}
\fxvdr{(1.583333,1) -- (2,1) -- (2,2) -- (1.583333,2) -- cycle}{(1.583333,1) -- (2,1) -- (2,2) -- (1.583333,2) -- cycle}{(1.583333,1) -- (2,1) -- (2,2) -- (1.583333,2) -- cycle}{0.42}
\fxvdr{(2.576831,0) -- (3,0) -- (3,1) -- (2.576831,1) -- cycle}{(3,0.55) -- (3,1) -- (2.576831,1) -- (2.576831,0.465366) -- cycle}{(2.7,1) -- (2.576831,1) -- (2.576831,0.589436) -- cycle}{0.42}
\fxvdr{(2.583333,1) -- (3,1) -- (3,2) -- (2.583333,2) -- cycle}{(2.583333,1) -- (3,1) -- (3,2) -- (2.583333,2) -- cycle}{(2.583333,1) -- (2.7,1) -- (3,2) -- (2.583333,2) -- cycle}{0.42}
\fxvdr{(3.575335,0) -- (4,0) -- (4,1) -- (3.575335,1) -- cycle}{(4,0.65) -- (4,1) -- (3.575335,1) -- (3.575335,0.565067) -- cycle}{}{0.42}
\fxvdr{(3.583333,1) -- (4,1) -- (4,2) -- (3.583333,2) -- cycle}{(3.583333,1) -- (4,1) -- (4,2) -- (3.583333,2) -- cycle}{}{0.42}
\fxvvel{1}{1}{0.5}{1.000}
\fxvvel{1}{1}{1.5}{1.000}
\fxvvel{1}{2}{0.5}{1.000}
\fxvvel{1}{2}{1.5}{1.000}
\fxvvel{1}{3}{0.5}{1.000}
\fxvvel{1}{3}{1.5}{1.000}
\fxvvel{1}{4}{0.5}{1.000}
\fxvvel{1}{4}{1.5}{1.000}
\fi
}%
  \begin{subfigure}[b]{0.33\textwidth}
    \centering
    \begin{tikzpicture}[x=1.2cm,y=1.2cm]
      \begin{scope}
        \clip (0,3) rectangle (3,6);
        \fxvpasses{1}
      \end{scope}
      \fxvarrows{1}{0.3}
      \draw[tr face] (3,4) -- (3,5);
      \node[tr small,anchor=west,fill=white,inner sep=1pt] at (3.12,4.5) {$f$};
      \node[tr small,fill=white,fill opacity=0.8,text opacity=1] at (2.35,4.5) {$\donor{f}$};
      \node[tr small,anchor=south,inner sep=1pt] at (2.88,6) {$\dr_f$};
      \draw[black!60,line width=0.3pt] (2.88,5.98) -- (2.92,5.6);
    \end{tikzpicture}
    \caption{Uniform faces.}
    \label{fig:vof:donating:uniform}
  \end{subfigure}%
  \hfill
  \begin{subfigure}[b]{0.29\textwidth}
    \centering
    \begin{tikzpicture}[x=1.7cm,y=1.7cm]
      \begin{scope}
        \clip (3,5) rectangle (5,7);
        \fxvpasses{2}
      \end{scope}
      \fxvarrows{2}{0.22}
      \draw[tr face] (4,6) -- (4,7);
      \node[tr vertex] at (4,6.5) {};
      \node[tr small,anchor=west,fill=white,inner sep=0.8pt] at (4.06,6.25) {$f_1$};
      \node[tr small,anchor=west,fill=white,inner sep=0.8pt] at (4.06,6.75) {$f_2$};
      \node[tr small,align=center,fill=white,fill opacity=0.8,text opacity=1]
        at (3.4,6.72) {$\donor{f_1}$\\$=\donor{f_2}$};
    \end{tikzpicture}
    \caption{Coarse donor, two segments.}
    \label{fig:vof:donating:coarse-fine}
  \end{subfigure}%
  \hfill
  \begin{subfigure}[b]{0.36\textwidth}
    \centering
    \begin{tikzpicture}[x=1.2cm,y=1.2cm]
      \begin{scope}
        \clip (1,0) rectangle (4,2);
        \fxvpasses{3}
      \end{scope}
      \fxvarrows{3}{0.3}
      \draw[tr face] (3,0) -- (3,1);
      \node[tr small,anchor=north west,inner sep=1pt] at (3.03,0.98) {$f$};
      \node[tr small,anchor=north,fill=white,inner sep=1pt] at (2.5,-0.03) {$\donor{f}$};
      \node[tr small,anchor=north,inner sep=1pt] at (3.6,-0.03) {$\aperture{f}=\chkVofCutAperture$};
    \end{tikzpicture}
    \caption{Cut donor and cut face.}
    \label{fig:vof:donating:cut}
  \end{subfigure}
  \caption{Donating regions~\cref{eq:vof:rectangle,eq:vof:depth,eq:vof:region} of
    the first sweep, along $\unitvec{1}$, at $\dt=\dtwy$. Dashed red: the
    rectangles $\drect{f}(\drdepth{f})$ of all faces with $\velo_f\ne0$, mapped
    to physical coordinates; arrows: $\velo_f$; hatched: the liquid part, whose
    volume is $\dt\meas{f}\volumeflux^{\liq}_f$; dots: cells with the frozen
    indicator $\indicator{\liq}_c=1$. (a,~b)~Shear flow
    $\velo=1+0.3\sin(\pi\coordk{2}/4)$ around a drop on a composite mesh
    (\code{vof_shear_fixture} of \code{figures/vof.jl}; unit coarse cells, the
    columns $4\le\coordk{1}\le8$ refined once); the depth is
    $\dt\abs{\velo_f}$ in units of the donor width, at most $\wycflval=\half$
    (\cref{eq:vof:cfl}), so the rectangles of the coarse cells are half as deep
    as those of the fine cells for equal velocities. (b)~Each of the two
    segments of the coarse side has its own rectangle over half of the coarse
    donor, with its own velocity and depth; open circle: the hanging vertex.
    (c)~Uniform flow $\velo=1$ over a sloped wall (hatched), PLIC lines of
    \code{:lhf_ghf}; the rectangle spans the whole face and its depth
    $\drdepth{f}=\chkVofCutDepth$ is set by the open area (shaded), whereas
    $\dt\abs{\velo_f}/\cellwidth{\donor{f}}{1}=\chkVofCutExtrusion$.}
  \label{fig:vof:donating}
\end{figure}
